\documentclass[12pt,a4paper]{article}

\usepackage[margin=1in]{geometry}
\usepackage[T1]{fontenc}
\usepackage{newtxtext}
\usepackage{setspace}
\usepackage{changepage}

\setlength{\parskip}{0.5\baselineskip}
\setlength{\parindent}{0pt}

\pagestyle{empty}
\onehalfspacing

\begin{document}

    \textbf{Quin Barnard}\\
    \textbf{27902277}\\
    \textbf{APP/1607896}\\
    \textbf{Bachelor of Arts Honours in General Linguistics: Research Proposal for 2027}\par

To whom it may concern,

I am proposing research in the field of Psycholinguistics, specifically the subfield Language Processing. I am interested in the effects of syntactic priming across different languages in a bilingual speaker of Afrikaans and English.\par

Afrikaans has been described as a negative concord language, yet this classification is at best reliant on context. Huddlestone and De Swart\footnote{Huddlestone, K. \& De Swart, H. 2014. A bidirectional Optimality Theoretic analysis of multiple negative indefinites in Afrikaans. \textit{Stellenbosch papers in Linguistics.} 43(1):137–164. DOI: 10.5774/43-0-167.} study the co-occurrence of Multiple Negative Indefinites (MNIs), such as \textit{“niemand, niks, nooit etc.”} producing a single negation meaning, that is to say a cancellation of the meaning that would exist were there no negation present. Classic prescriptive grammars of Afrikaans prohibit such interpretations, as MNIs are supposedly analogous to the Double Negative in English, yet it is prevalent in many modern varieties of Afrikaans.\\
It is well–known that a similar phenomenon occurs in some varieties of English, such as the utterance \textit{“We don't want nothing, not a thing from you”}\footnote{Szoke, J. 2010. Multiple Negation in English today. \textit{The Odd Yearbook} 8:139} producing the meaning \textit{“There is nothing we desire from you”} in, among other varieties, African American Vernacular English.\par

There is, however, no extensive literature on the effect of MNIs in South African English, and the literature on the same subject in Afrikaans is still developing as the different varieties of Afrikaans are not exhaustively mapped. Thus I propose the following research question:
\begin{adjustwidth}{1.5cm}{1.5cm}
    \begin{spacing}{1}
        \begin{center}
            \small
                \textit{Does the non–standard use of MNIs in either Afrikaans or South African English increase the likelihood of a single–negation interpretation in the other language among bilingual speakers?}
        \end{center}
    \end{spacing}
\end{adjustwidth}
The varieties of Afrikaans and English studied by such a question would have to be determined more precisely for the sake of rigour. As such a question would be studied in Stellenbosch, it would perhaps be easiest to consider the Cape Winelands varieties of Afrikaans and the English that bilingual speakers of this community speak. The study itself would consist of two groups, one with a target language of Afrikaans and an prime language of English, the other group being inverted. Participants would receive a short text containing MNI's with an unambiguous single–negation meaning in the prime language. They would then receive a similar text in the target language and be given a set of multiple choice questions on their interpretation of the latter. There would also need to be a control group for all variations of the texts to establish a baseline.\par

If this proposal is not suitable the department, I am more than willing to be assigned a topic, preferably within the field of Language Processing. Thank you for your consideration.
    
\end{document}
